% Short-edition Section 4. Description paragraphs reworded by the team on 2026-09-21 and condensed the same day. The \fr requirements match the full edition except that the business-rule cross-references (BR-n) are omitted, since the short edition has no business-rules section.
\section{System Features}\label{sec:features}

Each feature gives a short description with its priority, a stimulus/response pair, and the functional requirements that specify it. Identifiers match the full specification, which holds the complete requirement set; this section covers what is needed to understand and test Release~1. Stretch requirements are tagged [Stretch], have Low priority, and are built only if time allows.

% ------------------------------------------------------------------ 4.1
\subsection{Inventory and Resource Types}\label{feat:inventory}
\textbf{Description and priority.} The Inventory records every Resource, its Resource Type, and who holds it or has a Hold on it. An Administrator defines Resource Types as data, so a new kind of room, instrument, or pool needs no code change, and Resource Owners keep the Resources current (\FE{inventory}; \UC{define-type}, \UC{inventory}). \priority{High.} \textbf{Stimulus/response.} A Resource Owner reduces a pool's units; the system refuses any reduction below what is allocated or held.

\fr{Inventory.Type.Define}{The system shall let an Administrator define a Resource Type as data: a unique name, exactly one Allocation Class (Matched, Exclusive, or Metered), and named attributes of type text, number, date, or list of allowed values. Defining, changing, or retiring a type shall require no code or deployment change; a retired type accepts no new lines, a type referenced by any line is never deleted, and the built-in type Person (class Matched) cannot be retired.}
\fr{Inventory.Resource.Add}{The system shall let a Resource Owner add, edit, and retire an Exclusive or Metered Resource, filling in the type's attributes. A Metered Resource records the units it supplies; the Pool Capacity of a Metered type is the sum of the units of its unretired Resources. A retirement or reduction that would fall below the units allocated or held in any future hour shall be refused, naming the earliest such hour.}
\fr{Inventory.View.Holder}{The system shall show, for any Resource and window, whether it is free, on Hold, or allocated in each part of the window and by which Team and Request, and shall let a Member filter Resources by type, attribute value, and being free for a window. A Resource Owner has no approval step in any allocation.}

% ------------------------------------------------------------------ 4.2
\subsection{Member Profiles}\label{feat:profiles}
\textbf{Description and priority.} A Member's profile holds their Team, skills, and availability, and Matched lines are matched against it (\FE{profiles}). \priority{High.} \textbf{Stimulus/response.} A Member updates their skills or availability; the system rebuilds the Shortlist of any unsatisfied Matched line the Member now qualifies for.

\fr{Profile.Skills}{A Member shall belong to exactly one Team, set by an Administrator, and shall maintain their own skills, each chosen from the skill list an Administrator maintains, so that matching is exact.}
\fr{Profile.Availability}{A Member shall record availability as date ranges during which they may be lent; they are available for a window when every day of it lies inside a range. A Loan never changes these ranges, but a Person is in at most one Loan at a time: a Member with a Loan overlapping any day of a Matched line's window is not a Candidate for that line, and a Loan that ends early frees the Person from the day after its new last day. A Member sees their own profile and Loans; a Manager sees those of their Team's Members and may edit a Member's availability. A profile holds no personal data beyond name, email, Team, skills, and availability (\seref{data}).}

% ------------------------------------------------------------------ 4.3
\subsection{Suggestions}\label{feat:suggestions}
\textbf{Description and priority.} A Suggestion is one Request Line that a Member proposes to their own Manager; it affects nothing until the Manager accepts it into a Request and submits. Declined Suggestions are kept as the data for the Future unmet-demand analytics (\FE{suggestions}; \UC{suggest}). \priority{Medium.} \textbf{Stimulus/response.} A Member submits a Suggestion; the Manager accepts it into a Draft Request or declines it with an optional reason, and the Member is notified.

\fr{Suggest.Create}{The system shall let a Member create a Suggestion addressed to their own Manager containing exactly one proposed line (Resource Type, quantity, window, for a Matched line a skill set or a named Person, optional note). A Suggestion is Open, Accepted, or Declined; only Open ones change state, none is ever deleted, and a Suggestion has no effect on the Inventory, Holds, or the Ledger until the line it was accepted into is submitted.}
\fr{Suggest.Accept.Existing}{The system shall let the Manager accept an Open Suggestion into a Draft Request as a new line or into an existing line of the same Resource Type, so that several Suggestions point at one line; an Accepted Suggestion records the line it became and keeps its original ask unchanged. The Manager may instead decline it with an optional reason, and the Member is notified.}

% ------------------------------------------------------------------ 4.4
\subsection{Draft and Submit a Request}\label{feat:draft}
\textbf{Description and priority.} A Manager builds a Request in Draft, marks each line Required or Optional, sees a Quote per line, and submits; submission is refused while the Required lines exceed the Available Balance (\FE{draft}; \UC{draft}). \priority{High.} \textbf{Stimulus/response.} The Manager submits; the system locks every Quote and evaluates the Request (Section~\ref{feat:lifecycle}), or shows the shortfall and leaves it in Draft.

\fr{Draft.Create}{Only a Manager shall create, edit, submit, or withdraw a Request, and only on behalf of their own Team. Withdrawing from Draft, Pending, or Waiting releases all Holds and held credits; editing a Waiting Request returns it to Draft, removes its lines from waitlists, and recomputes Quotes on resubmission.}
\fr{Draft.Line.Add}{A Request Line shall name a Resource Type, a quantity, and a window; its Allocation Class is taken from the type. A Matched line gives a skill set and quantity or a named Person with quantity one; an Exclusive line may name a specific Resource, a weekly recurrence, and the Priority mark; a Metered line gives units held concurrently. A line whose window does not start in the future, or whose type is retired, shall be refused.}
\fr{Draft.Line.Required}{Every line shall be Required (default) or Optional. The Manager may split a line into two of the same type to mark a part Optional; there is no per-line partial-acceptance setting.}
\fr{Draft.Quote.Line}{While drafting, the system shall show each Exclusive and Metered line's Quote computed as in \FR{Ledger.Price.Formula} at the current settings, the sums for Required and Optional lines, the Available Balance, and any shortfall; Matched lines show no Quote because they cost nothing. The system shall refuse submission while the Required sum exceeds the Available Balance.}
\fr{Draft.Submit.Lock}{On submission the system shall lock every Quote, record the submission time (used for waitlist order), and evaluate every line as specified in Section~\ref{feat:lifecycle}; a locked Quote never changes and the Charge posted equals it.}

% ------------------------------------------------------------------ 4.5
\subsection{Matched Allocation}\label{feat:matched}
\textbf{Description and priority.} A Matched line asks for people by skills and quantity, or by name. The system builds a Shortlist that leaves out anyone already on a Loan in the window and lists the least-lent first; the requesting Manager picks Candidates, and each Candidate's own Manager approves within the Objection Window or the Capacity Planner decides. An approved Candidate becomes a Loan when the Request is Allocated; the borrowing Team is charged nothing and the lending Team earns the Day Rate (\FE{matched}; \UC{shortlist}, \UC{approve}, \UC{escalate}). \priority{High.} \textbf{Stimulus/response.} A lending Manager approves or refuses; the last Required approval allocates the Request, and a refusal releases every Hold and asks the requesting Manager for a replacement.

\fr{Matched.Shortlist.Build}{When a Request enters Pending, the system shall build for each skill-based Matched line a Shortlist of every Member outside the requesting Team who holds every listed skill, is available for the whole window, and has no Loan overlapping any day of it, ordered by fewest person-days lent in the current Period (the days of every formed Loan that fall in the Period), ties by name; a named Person is checked only for availability and an overlapping Loan. The requesting Manager shall select exactly the line's quantity of Candidates; the system never selects automatically, and a Shortlist shorter than the quantity leaves a Required line unsatisfied and the Request Waiting.}
\fr{Matched.Approve.Route}{For each selected Candidate the system shall send one approval request to that Candidate's Manager, stating the requesting Team, the window, and the Earning; the lending Manager approves or refuses with an optional reason, recorded with actor and timestamp. No approval request is sent to, or accepted from, the requesting Manager for their own Team.}
\fr{Matched.Refuse.Replace}{When an approval is refused, the system shall ask the requesting Manager to select a replacement; the Objection Window shall not restart. If no Candidate remains, or the line named a Person, a Required line places the Request in Refused and an Optional line is dropped.}
\fr{Matched.Escalate}{When the Objection Window lapses with an approval outstanding, the system shall present the outstanding approvals to the Capacity Planner, who approves or refuses on the silent Manager's behalf; every later approval on that Request goes to the Capacity Planner directly.}
\fr{Matched.Loan.Form}{When the Request becomes Allocated, each approved Candidate shall become a Loan for the line's window and an Earning of the Day Rate times the person-days shall post to the lending Team; an Optional line approved later forms its Loan then. The Candidate is notified when selected and the Person when the Loan forms; the Person cannot decline in Release~1.}

% ------------------------------------------------------------------ 4.6
\subsection{Exclusive Allocation}\label{feat:exclusive}
\textbf{Description and priority.} An Exclusive line asks for a Resource held whole over a window. The system checks allocations and Holds for conflicts, picks a free Resource when none is named, and waitlists an unmet Required line. A Priority line carries a premium, can Bump a standard line outside the bump cutoff, and is never Bumped itself (\FE{exclusive}; \UC{bump}, \UC{waitlist}). \priority{High; Priority and Bump are Medium.} \textbf{Stimulus/response.} A Priority line finds no free Resource; the system Bumps a standard line, refunds it in full, and moves it to the head of the waitlist.

\fr{Exclusive.Conflict}{An instance shall be free for a window when no allocation and no Hold on it overlaps any part of the window; a line is satisfiable when its quantity of free instances exists for every occurrence on the same instance. When no Resource is named the system picks the free instance with the fewest allocated hours in the Period.}
\fr{Exclusive.Waitlist.Join}{A Required Exclusive line that is unsatisfiable shall join the waitlist for its Resource Type and window; an Optional line never does. Waitlists are ordered Priority lines first, then by submission time, with a Bumped line at the head; a waiting line leaves the waitlist when its window start passes and its Request stays Waiting; the Manager sees each waiting line's position. When capacity frees through a Release, a Bump, or an expired Hold, the system offers it to the waiting lines in order and satisfies the first whose window fits, under \FR{Lifecycle.Waiting.Reevaluate}.}
\fr{Exclusive.Priority.Bump}{The Manager shall be able to mark an Exclusive line as a Priority line while drafting, at the priority premium; Priority is offered for Exclusive lines only in Release~1. When a Priority line finds no free instance, the system shall Bump the standard lines, allocated or on Hold, that block the instance needing the fewest Bumps, preferring the latest-submitted, provided no Bumped window or occurrence starts within the bump cutoff (example default 24~hours); a Priority line that cannot Bump joins the waitlist ahead of all standard lines. A Priority line is never Bumped and never Bumps another Priority line.}
\fr{Exclusive.Bump.Effect}{When a line is Bumped the system shall, in one transaction, post a Refund equal to its full Charge, place it at the head of its waitlist, notify its Manager with the reason, move its Request to Waiting if the line was Required in an Allocated Request while the other lines stay allocated, and record the Bump with both lines, actor, and timestamp.}

% ------------------------------------------------------------------ 4.7
\subsection{Metered Allocation}\label{feat:metered}
\textbf{Description and priority.} A Metered line asks for units held concurrently over a window. It fits when, in every hour, allocated plus held plus requested units stay within the Pool Capacity and within the Team's Quota where one is set; the Quota stops one Team with many credits from taking the whole pool. Service tiers that Bump and requested-versus-used tracking are Stretch (\FE{metered}). \priority{High.} \textbf{Stimulus/response.} A Metered line is submitted; the system checks it hour by hour and holds, allocates, or waitlists it.

\fr{Metered.Line.Form}{A Metered line shall name a Resource Type, units to hold concurrently, and a window; no specific Resource and no Priority mark is offered; the line is satisfied with its full units for its full window or not at all. Its Quote is the sum over hours of units times Unit Price times Price Band multiplier.}
\fr{Metered.Pool.Check}{A Metered line shall be satisfiable only if, for every hour of its window, the units allocated plus held for that type plus the line's units do not exceed the Pool Capacity and, where the Capacity Planner has set a Quota for the requesting Team on the type, the Team's allocated plus held plus requested units do not exceed it. A Required line that does not fit joins the waitlist under the rules of \FR{Exclusive.Waitlist.Join} and is satisfied when its units fit in every hour.}
\fr[Stretch]{Metered.Tier.Define}{The system shall let the Capacity Planner define service tiers for a Metered Resource Type, each with a premium multiple on the Unit Price, so that a Manager can mark a Metered line with a tier in the same way as a Priority line.}
\fr[Stretch]{Metered.Tier.Bump}{A line of a higher tier shall be able to Bump units of lower-tier lines under the rules of \FR{Exclusive.Priority.Bump} to \FR{Exclusive.Bump.Effect}, the bump cutoff included; the Bumped line shall receive a full Refund and the head of the waitlist.}
\fr[Stretch]{Metered.Usage.Record}{The system shall let a Resource Owner record, at the end of a Period, the units actually consumed by each Metered allocation of that Period.}
\fr[Stretch]{Metered.Usage.Compare}{The dashboard shall show requested versus used units per Team and Resource Type per Period. Recorded usage shall have no effect on the Ledger.}

% ------------------------------------------------------------------ 4.8
\subsection{Request Lifecycle and Holds}\label{feat:lifecycle}
\textbf{Description and priority.} Exclusive and Metered lines resolve at once; Matched lines wait for approval. Allocating the room before the approvals arrive charges the Team for a room it may never use, so a Request that enters Pending places a Hold on every Exclusive and Metered line that fits. A Hold blocks capacity, posts no Charge, and never outlives the Objection Window (\figref{fig:states}; \FE{lifecycle}). \priority{High.} \textbf{Stimulus/response.} The last Required approval arrives; the system commits every Hold and posts every Charge and Earning in one transaction. An approval is refused or the window lapses; the Holds release and the Request moves to Waiting.

\fr{Lifecycle.States}{A Request shall be in exactly one of Draft, Submitted, Pending, Allocated, Waiting, Withdrawn, or Refused, with exactly the transitions of \figref{fig:states}; every state change of a Request or line records actor, change, and timestamp.}
\fr{Lifecycle.Evaluate}{On submission the system shall, in one transaction, decide for every Exclusive and Metered line whether it is satisfiable now and for every Matched line whether an approval is needed, and move the Request to Allocated if no approval is needed and every Required line is satisfiable, to Pending if any line needs approval, and otherwise to Waiting.}
\fr{Lifecycle.Hold.Place}{When a Request enters Pending the system shall place a Hold on every satisfiable Exclusive and Metered line. A Hold blocks other Requests exactly as an allocation does, posts no Charge, and its locked Quote counts as held credits so that concurrent Pending Requests cannot together exceed the Balance. A Hold commits when the Request becomes Allocated and otherwise releases when the Request leaves Pending or the Objection Window lapses; none is placed after the window has lapsed.}
\fr{Lifecycle.Window.Start}{The Objection Window shall start when the Request first enters Pending, last the length set by the Capacity Planner (example default two working days), and not restart when a Candidate is replaced. On lapse with a Required approval outstanding, the system shall release every Hold, move the Request to Waiting, and escalate, within five minutes (\peref{jobs}).}
\fr{Lifecycle.Allocate}{When every Required line is satisfied, the system shall in one transaction commit every Hold, post a Charge equal to each locked Quote, form the Loans and post the Earnings, drop Optional lines that are unsatisfiable, unapproved, or no longer affordable (telling the Manager which and why), move the Request to Allocated, and notify the Manager; every other state change, Bump, waitlist fill, escalation, and approval request is notified in-app.}
\fr{Lifecycle.Waiting.Reevaluate}{Whenever capacity a Waiting Request waits for frees, or an approval on it arrives, the system shall check whether every Required line can be satisfied at that moment; if so it allocates the Request at the locked Quotes, otherwise the freed capacity passes to the next waiting line and the Request's lines keep their positions. A replacement Candidate picked while the window still runs returns the Request to Pending with Holds re-placed.}

% ------------------------------------------------------------------ 4.9
\subsection{Release and Refund}\label{feat:release}
\textbf{Description and priority.} A Manager can Release all or part of an allocated Exclusive or Metered line for a Refund set by the Resource Type's Refund Schedule, a step function of lead time that never reaches 100\%: holding capacity is never free, but an early Release beats a silent no-show. Ending a Loan early prorates the lending Team's Earning (\FE{release}; \UC{release}). \priority{High.} \textbf{Stimulus/response.} The Manager Releases a line; the system posts the Refund, frees the capacity, and offers it to the waitlist.

\fr{Release.Line}{The system shall let the owning Team's Manager Release an allocated Exclusive or Metered line in whole or in part before it is consumed; a partial Release shortens the window, removes occurrences, or reduces units, the rest keeps its allocation, instance, and prices, and the freed capacity is offered to the waitlist at once (\FR{Exclusive.Waitlist.Join}). Capacity neither used nor Released is never freed automatically and appears on the dashboard as idle.}
\fr{Release.Refund}{On Release the system shall post one Refund equal to the Charge for the released part times the fraction from the Refund Schedule for the lead time to that part's start (example default 75\% at seven or more days, 25\% under seven days, 0\% once started); no step may be 100\% or more, and the Refund returns to the parts of the Account the Charge was drawn from.}
\fr{Release.Loan.Prorate}{The borrowing Manager shall be able to end a Loan early; the system shall post to the lending Team an Earning correction so that its Earning equals the Day Rate times the days actually lent. No Refund Schedule applies. This correction is the only Transaction that may take a Balance below zero, and until a Grant or Earning restores it the Team can submit no Request with a Required Exclusive or Metered line.}

% ------------------------------------------------------------------ 4.10
\subsection{Ledger and Pricing}\label{feat:ledger}
\textbf{Description and priority.} The Ledger holds each Team's Account and every Transaction. Exclusive and Metered lines are charged credits, Matched lines nothing, and lending earns the Day Rate; granted credits expire with the Period while earned credits carry over. Prices vary by Price Band and lock into the Quote at submission; a demand surcharge is Stretch (\FE{ledger}; \UC{period}, \UC{dashboard}). \priority{High.} \textbf{Stimulus/response.} A Period begins; the system posts the Grants and expires the previous Period's granted credits.

\fr{Ledger.Account}{Each Team shall have one Account holding granted, earned, and held credits; Balance is granted plus earned, Available Balance is Balance minus held, and held credits equal the locked Quotes of the Team's lines on Hold. For each Team and Period the system shall produce a Statement with opening granted and earned credits, every Transaction, the amount expired, and closing balances; a Manager sees their own Team's, the Capacity Planner every Team's.}
\fr{Ledger.Transaction.Kinds}{Every Transaction shall be a Grant, Charge, Earning, or Refund, recording Team, amount, Period, timestamp, actor or job, and the Request Line concerned; Transactions are never updated or deleted, and a correction is a new Transaction referencing the one it corrects.}
\fr{Ledger.Grant.Period}{At the start of each Period the system shall post every Team's Grant as set by the Capacity Planner, who may also post a mid-Period Grant with a reason; at the end of a Period it shall set granted credits to zero, record the expired amount on the Statement, and leave earned credits untouched, within fifteen minutes of the boundary (\peref{jobs}). A Charge draws granted credits first, then earned, posts only when the Request becomes Allocated, equals the locked Quote, and never takes a Balance below zero.}
\fr{Ledger.Price.Formula}{The Capacity Planner shall set a Unit Price per instance-hour or unit-hour for each Exclusive and Metered type and non-overlapping Price Bands (weekdays, time of day, multiplier; unbanded hours 1.0; example defaults weekdays 09:00--17:00 $\times$1.5, other hours $\times$0.5). A line's Quote is the sum over each hour of each occurrence of quantity times Unit Price times band multiplier, times the priority premium (example default 2$\times$) for a Priority line.}
\fr{Ledger.Earning.Post}{An Earning shall post to the lending Team's earned credits when a Loan forms, equal to the Day Rate times the Loan's person-days, a person-day being a Monday-to-Friday calendar day inside the window; the Day Rate is one value for every Person.}
\fr[Stretch]{Ledger.Surcharge}{The Capacity Planner shall be able to enable a demand surcharge per Resource Type: an additional multiplier on the Unit Price that rises with the fraction of the type's capacity already allocated or held in the requested window at quote time, applied in \FR{Ledger.Price.Formula} and locked with the Quote.}

% ------------------------------------------------------------------ 4.11
\subsection{Showback Dashboard}\label{feat:dashboard}
\textbf{Description and priority.} The dashboard is showback: utilization and idle capacity per Resource Type, credits earned and spent per Team, and each Period's Statement (\FE{dashboard}; \UC{dashboard}). \priority{Medium.} \textbf{Stimulus/response.} A Member opens the dashboard; the system shows the current Period's figures for their role, or any past Period they pick.

\fr{Dashboard.Utilization}{For each Exclusive and Metered Resource Type and Period the dashboard shall show utilization as a percentage (allocated instance-hours over available, or allocated unit-hours over Pool Capacity times hours), idle capacity (available minus allocated) for the Period to date and the coming seven days, and the number of lines waiting.}
\fr{Dashboard.Credits}{For each Team and Period the dashboard shall show Grants, Charges, Earnings, Refunds, credits expired, Balance and Available Balance, and person-days lent and borrowed. A Manager sees company-wide utilization and only their own Team's credits and Loans; the Capacity Planner every Team; a Resource Owner the Resources they own. Figures reflect every allocation, Release, and Transaction completed before the page was requested.}
\fr[Stretch]{Dashboard.Usage}{The dashboard shall show requested versus used units per Team and Metered Resource Type when usage has been recorded (\FR{Metered.Usage.Compare}).}

% ------------------------------------------------------------------ 4.12
\subsection{Administration}\label{feat:admin}
\textbf{Description and priority.} An Administrator creates Members and Teams, grants roles, maintains the skill list, and defines Resource Types, and sets no policy; every Capacity Planner setting is specified here so that each has one home (\FE{admin}; \UC{define-type}, \UC{period}). \priority{High.} \textbf{Stimulus/response.} The Capacity Planner changes a setting; it applies to future Quotes, Grants, and windows and never alters a locked Quote or a posted Transaction.

\fr{Admin.Member.Create}{An Administrator shall be able to create and deactivate Members (name, email, username, initial password, Team), create, rename, and retire Teams, move a Member between Teams, maintain the skill list, and grant and revoke the roles Manager, Resource Owner, Capacity Planner, and Administrator; each Team has exactly one Manager at any time, and a deactivated Member cannot log in and appears on no Shortlist. Users authenticate by username and password (\seref{password}); roles are enforced on the server; Administrator screens contain no policy setting; every setting change is recorded with actor, old and new value, and timestamp.}
\fr{Planner.Settings}{The Capacity Planner shall be able to set as data, without deployment: the Period (default one calendar month); each Team's Grant per Period; each Exclusive and Metered type's Unit Price, Price Bands, and Refund Schedule; the Day Rate; any Team's Quota on a Metered type; the priority premium; each Exclusive type's bump cutoff; and the Objection Window. A changed price affects only Quotes not yet locked; a changed window, cutoff, or schedule affects events after the change. The Capacity Planner shall see the escalated Requests and decide each outstanding approval from that list.}
